\ifdefined\MyWebBook\else
\documentclass[../textbook.tex]{subfiles}
\begin{document}
\fi
\bookchapter{训练工程优化}{ch:efficient-finetuning}

微调的工程优化以任务质量要求为约束，降低显存需求、计算量和完成时间。监督微调规定数据与损失，参数高效微调规定可训练参数集合；在此基础上，还需选择数值精度、批量组织与并行策略。本章讨论这些方法在单设备、单机多卡和多机执行中的组合。LoRA 的参数化与梯度机制见\bookxref{ch:parameter-efficient-finetuning}，跨设备通信和分片机制见\bookxref{ch:distributed-training}。

比较方案时，应固定监督词元、验证集和质量要求，同时记录各设备峰值显存、全局有效监督词元吞吐、达到目标质量的时间与设备时成本。

\section{微调显存预算}

设冻结参数数量为 $P_f$，可训练参数数量为 $P_t$，每个冻结参数的实际存储开销为 $b_f$，每个可训练参数及其梯度、优化器状态的总字节数为 $b_t$。单设备峰值可表示为
\begin{equation}
 M_{\mathrm{peak}}=P_fb_f+P_tb_t+M_{\mathrm{act}}+M_{\mathrm{temp}}+M_{\mathrm{runtime}}.
\end{equation}
量化尺度、分组元数据和对齐开销应计入 $b_f$；主参数副本与优化器矩状态应计入 $b_t$。激活项取决于序列长度、微批量、可训练模块的位置和重计算策略。冻结参数仍然参与前向计算；当梯度需要经过冻结层传向更早的适配器时，该层的反向计算也必须保留。

\begin{figure}[htbp]
  \centering
  \includegraphics[width=0.92\textwidth]{\subfix{../../figures/theory/peft-memory-breakdown.pdf}}
  \caption{微调各范式显存占用结构对比}
  \label{fig:peft-memory-breakdown}
\end{figure}

\begin{example}[微调的完整显存预算]
例如冻结基座有 $10^9$ 个参数，每参数实际开销为0.6字节，可训练参数有 $10^7$ 个，每参数训练状态为16字节，两项合计 $7.6\times10^8$ 字节。若激活、临时空间和运行时另占 $3\times10^9$ 字节，则峰值为 $3.76\times10^9$ 字节。只报告可训练参数比例会遗漏其中占比最大的一项。实际容量可通过完整训练步的峰值测量校正。
\end{example}

\begin{table}[htbp]
  \centering
  \small
  \setlength{\tabcolsep}{3.5pt}
  \caption{全参数微调与 PEFT 预算对比}
  \label{tab:peft-comparison}
  \begin{tabularx}{\linewidth}{lccccX}
    \toprule
    微调范式 & 基座精度 & 适配器 & 可训参数比 & 显存瓶颈项 & 典型适用场景 \\
    \midrule
    全量微调 (FFT) & BF16 & -- & $100\%$ & 优化器状态与梯度 & 领域差异显著的持续预训练与深度对齐 \\
    LoRA & BF16 & BF16 & $0.05\%\sim 1\%$ & 冻结权重与激活 & 中等算力指令微调与多租户适配器路由 \\
    QLoRA & 4-bit NF4 & BF16 & $0.05\%\sim 1\%$ & 动态激活张量 & 消费级显卡或单卡严苛预算下的模型适配 \\
    \bottomrule
  \end{tabularx}
\end{table}

\section{梯度累积的归一化}

设一次更新包含 $G$ 个微批次，第 $j$ 个微批次的有效监督词元数为 $n_j$、损失和为 $S_j$。在固定标签掩码下，令 $N=\sum_j n_j>0$，则
\begin{equation}
 \begin{aligned}
  \mathcal L &= \frac{\sum_j S_j}{N}, \\
  \nabla\mathcal L &= \sum_{j:n_j>0}\frac{n_j}{N}\nabla\overline{\mathcal L}_j, \quad \overline{\mathcal L}_j=\frac{S_j}{n_j}.
 \end{aligned}
\end{equation}
逐微批次反向传播前因此需要按有效监督词元数加权。提示词元和填充词元可能参与前向计算，却不属于损失分母。若两个微批次分别有2和6个监督词元，平均梯度分别为1和3，则目标梯度为 $\frac{2\times1+6\times3}{8}=2.5$，简单平均得到2，改变了优化目标。

\begin{warning}[梯度累积的微批次加权归一化陷阱]
在变长序列或包含提示掩码的指令微调中，各个微批次包含的有效监督词元数 $n_j$ 往往不相等。若在每个微批次内部独立求均值损失并在更新时除以累积步数 $G$，将对短回答或监督稀疏批次赋予过高权重，隐式篡改真实优化目标。正确工程实现应保持微批次损失总和，并在全局规约时严格按总有效词元数 $N$ 进行归一化。
\end{warning}

所有微批次应使用同一参数版本；一次累积窗口结束后再归一化、裁剪、更新参数和优化器状态。学习率日程按实际更新计时。累积降低单次前向的容量需求，但不减少每个样本所需的计算。涉及批内统计的层、随机掩码以及浮点归约顺序会影响累积与完整批次的数值一致性。

\section{训练批次组织}

对批内长度 $\ell_1,\ldots,\ell_B$，统一填充到 $\ell_{\max}$ 时，填充位置比例为
\begin{equation}
 \rho_{\mathrm{pad}}=1-\frac{\sum_j\ell_j}{B\ell_{\max}}.
\end{equation}
长度为 $(2,2,6,6)$ 的样本合为一个批次需要24个位置，分为两个等长批次需要16个位置。长度分桶还需保留随机化和任务混合，避免形成固定的长度课程。注意力交互和前馈计算对序列长度具有不同的依赖关系，端到端加速比需结合各算子的执行时间及调度开销测量。

将多条指令装入同一序列时，应维护样本边界、位置语义和回答区域掩码。独立样本通过注意力掩码阻止跨样本读取，并用终止词元标记回答结束。截断则需要明确保留哪些监督区域，并检查是否系统性地排除了长回答或特定任务。

\section{训练内存优化}

混合精度分别规定权重存储、乘法输入、累加、梯度和优化器状态的格式。损失缩放后的梯度须在裁剪前还原；出现非有限值时，参数更新和优化器时间应遵循一致的跳步规则。数值稳定性的机制与恢复条件见\bookxref{ch:training-stability}。

激活重计算保存部分前向边界，在反向时重新求值区间内部。对 $L$ 层均匀链式网络，每层一个大小为 $A$ 的激活单位，每隔 $k$ 层保存边界时，近似存储为
\begin{equation}
 M_{\mathrm{act}}(k)\approx A\left(\frac{L}{k}+k\right).
\end{equation}
将 $k$ 暂视为正实数，导数在 $k=\sqrt L$ 时为零，得到近似最小值 $2A\sqrt L$。若 $L=64$、$A=16$ MiB，选择 $k=8$ 得到256 MiB，全部保存则为1024 MiB。边界、临时张量和非均匀层会改变实际峰值。

\begin{table}[htbp]
  \centering
  \caption{激活重计算策略对比与开销分析}
  \label{tab:activation-checkpointing}
  \begin{tabularx}{\textwidth}{lp{28mm}p{30mm}X}
    \toprule
    重计算策略 & 显存开销 ($M_{\mathrm{act}}$) & 计算开销 (FLOPs) & 适用场景与工程权衡 \\
    \midrule
    无重计算 (None) & $\mathcal{O}(L \cdot s \cdot d)$ & $3LF$（标准前反向） & 短上下文（如 $s \le 2048$）、显存充足且追求极限吞吐 \\
    全量重计算 (Full) & $\mathcal{O}(s \cdot d)$ & $4LF$（约增加 $33\%$ 开销） & 显存极度紧张，通过增加计算时间换取超长序列或大批量 \\
    选择性重计算 (Selective) & $\mathcal{O}(\frac{L \cdot s \cdot d}{r})$ & 约 $3.1\sim 3.3LF$（开销增加 $<10\%$） & 仅重算显存占比大但计算量小的注意力中间激活（如 Softmax/Dropout） \\
    块级均分重计算 ($\sqrt{L}$) & $\mathcal{O}(\sqrt{L} \cdot s \cdot d)$ & 约 $4LF$ & 经典理论最优子线性显存策略，链式 Transformer 层分段边界存储 \\
    \bottomrule
  \end{tabularx}
\end{table}

重计算应复用相同的随机函数状态，并避免重复修改持久缓冲。若每层前向代价为 $F$，反向为 $2F$，额外重算一次前向会把近似工作量从 $3LF$ 增加到 $4LF$。节省的显存可用于更长序列或更大微批量，总时间是否下降则要测量完整训练步才能确定。相关存储构造见\autocite{chen2016sublinear}。

\section{分布式微调方案的选用}

先区分\emph{容量瓶颈}与\emph{吞吐瓶颈}。完整微调状态能够被单设备完全容纳时，可使用分布式数据并行（DDP）让不同副本处理不同数据分片以提升全局吞吐；DDP 要求每个副本保留完整的模型权重与优化器状态，无法突破单设备物理显存上限。参数高效微调虽然大幅降低需要跨卡全规约（All-Reduce）同步的梯度体积，但冻结基座的显存占用、前向计算以及必要的反向传播路径依然存在。

\keyconcept{微调工程优化必须明确区分容量瓶颈与吞吐瓶颈：全参数微调受限于优化器状态与参数冗余存储，应优先采用 ZeRO 或 FSDP 分片；参数高效微调虽压缩了可训参数开销，但在长上下文下极易受制于动态激活显存，必须协同采用选择性激活重计算或序列并行。}

若瓶颈是优化器状态、梯度或参数的重复存储，应根据占比选择 ZeRO 或全分片数据并行（FSDP）等状态分片方案。可训练参数很少时，优化器状态分片的绝对收益可能较小；基座参数分片能够降低常驻容量，却增加参数收集及临时缓冲。若显存瓶颈主要来自超长序列带来的动态激活张量，应优先采用选择性激活重计算、缩减微批次或引入序列并行（Context Parallelism），仅扩大参数分片组无法缓解各层前向激活张量的单卡驻留压力。张量或流水并行可进一步切分模型执行，但需要评估通信、流水气泡和实际设备拓扑。这些机制的计算与通信推导集中于\bookxref{ch:distributed-training}。

全局有效批量取决于数据并行副本数、微批量、累积次数和有效监督词元数。增加设备后，应明确保持全局批量还是增加每步监督量；两者改变更新次数和学习率日程的方式不同。张量并行成员处理同一批数据时，全局监督词元只应计入一次。各副本监督长度不等时，还需按全局词元数归一化，保证执行策略保持同一训练目标。

\begin{example}[墙钟时间与设备时成本]
例如，固定监督预算和质量要求下，单设备耗时100分钟，两设备耗时60分钟，则时间加速比为 $\frac{5}{3}$，设备时从100设备分钟增至120设备分钟。缩短端到端墙钟耗时（提高时效性）与降低总设备时（控制算力成本）存在工程权衡：并行度的提升往往伴随跨节点通信开销与梯度同步等待，导致并行效率下降并抬高总计算开销。方案对比必须建立在相同有效监督词元预算与收敛质量的基础上，多机执行还需综合考量网络通信尾部延迟、分布式检查点写入与数据加载瓶颈。
\end{example}

量化基座、适配器训练与状态分片的组合，需要核对存储格式、可训练与冻结参数的组织、反向传播、优化器和检查点恢复是否相容。组合方案需验证完整更新和恢复过程中的数值一致性与资源需求。

\section{微调执行评估}

先固定模型、分词器、模板、回答掩码和可训练模块，再用短序列检查前向损失与梯度。随后按目标长度测量一个完整更新窗口，计入优化器首次分配的状态；依据峰值调整微批量、累积步数和重计算范围。每次调整后同时检查损失分母、有效批量和随机状态，保证资源设置与训练目标一致。

比较全参数微调、LoRA 和 QLoRA 及其并行组合时，记录验证质量、全局有效监督词元吞吐、各设备峰值显存、总更新次数、墙钟时间和设备时。分布式执行还应核对数据分片、梯度同步、全局归一化、非有限值跳步与恢复一致性。检查点应包含适配器或完整参数、基座身份、优化器状态、学习率进度、随机状态与数据游标；分片检查点还需明确恢复所需的并行配置或重分片能力。发布时核对适配器合并、量化转换与服务模板造成的质量变化。

\chapterreferences
\myendchapterdocument
